\documentclass[10pt]{article}

\usepackage[margin=0.6in]{geometry}
\usepackage[T1]{fontenc}
\usepackage{lmodern}
\usepackage{microtype}
\usepackage{booktabs}
\usepackage{tabularx}
\usepackage{enumitem}

\pagestyle{empty}
\setlength{\parindent}{0pt}
\setlength{\parskip}{2pt}
\setlist{leftmargin=1.1em, itemsep=0.5pt, parsep=0pt, topsep=1pt}
\newcolumntype{Y}{>{\raggedright\arraybackslash}X}
\newcommand{\head}[1]{\par\vspace{4pt}{\large\bfseries #1}\par\vspace{1pt}}
\newcommand{\person}[2]{\par\vspace{2pt}\textbf{#1} \textit{(planned: #2)}}

\begin{document}

\begin{center}
  {\large\bfseries ECE473 Milestone II Update: New Rain Insurance Payouts and
  the AgriHedge App}\\[1pt]
  Maya, Evan, and Jack \quad\textbullet\quad November 12, 2026
\end{center}
\vspace{-8pt}

\textbf{Summary.} Milestone I showed that today's Kalshi rain contracts copy
USDA's rainfall insurance (PRF) poorly and at a higher cost. In Milestone II we
found out why, used that to design new kinds of insurance payouts that Kalshi
contracts can build exactly, and built an app that prices them for any farm.
Below, \emph{fit} is the share of a farm's year-to-year insurance payout swings
that a hedge reproduces, measured on years it was not trained on.

% ---------------------------------------------------------------------------
\head{1. Planned vs.\ actual progress}

\person{Jack}{find the cheapest Kalshi portfolio that matches a farm's PRF
payout; live pricing for the app}
\begin{itemize}
  \item \textbf{W4:} built an optimizer that picks contracts while accounting
    for fees, cash tied up until settlement, and how many contracts are
    actually for sale. Even the best portfolio reaches a fit of only about one
    third, so a better optimizer cannot fix the problem on its own.
  \item \textbf{W5:} built the backend that turns a farmer's plan into an
    order at live Kalshi prices, so every quote is one a farmer could pay
    today.
\end{itemize}

\person{Maya}{explain where the remaining hedge error comes from; farmer
screen}
\begin{itemize}
  \item \textbf{W4:} added real-world limits one at a time to measure each.
    Geography costs the most: Kalshi settles on an airport gauge 30 to 55
    miles from the farm. Monthly contracts (PRF insures two-month windows)
    and thin order books come next. This separates what better contracts
    could fix from what they cannot.
  \item \textbf{W5:} built the farmer screen: click your land, see the nearest
    station, how well it tracks your farm, and the cost. It showed the best
    use is the months PRF leaves uncovered (fit 0.61 at 9 of 11 farms, versus
    0.23 for copying PRF's own months).
\end{itemize}

\person{Evan}{find the risks today's contracts cannot cover and propose fixes;
test the app against PRF}
\begin{itemize}
  \item \textbf{W4:} distance to the airport gauge decides whether a hedge
    works. Across 150 random locations, fit is 0.47 within 25 miles of a
    Kalshi station but 0.15 beyond 100 miles, and only 7.4\% of US PRF land
    lies within 50 miles of one. Summer storm rain is too local for any
    airport gauge, so only 3 of the 36 two-month windows our farms insure can
    be copied well.
  \item \textbf{W5:} researched exotic payouts, such as exponential ones that
    pay little until a deep drought. Kalshi's one-inch strikes are too coarse
    for them (63\% of variance reproduced); half-inch strikes would give 86\%.
\end{itemize}

\textbf{Beyond the plan (shared equally):} new payouts and the advisor.
\begin{itemize}
  \item \textbf{Payouts built to order.} Kalshi lists a contract at each inch
    of monthly rain, so combining them can pay any amount at each rain level.
    We designed insurance that pays little in a mild dry spell and a lot in a
    severe drought, instead of PRF's straight-line payout. Since it is built
    from listed contracts, the insurer takes on no rain risk.
  \item \textbf{Why drought-weighted.} The best contract for a risk-averse
    farmer paying Kalshi's roughly 20\% markup puts almost all of its payout
    into the driest months. Our drought-weighted payouts get 90 to 93\% of its
    value; PRF's straight-line shape gets 67\%. A \$6,000 November contract at
    Dallas costs \$1,205 and pays \$3,727 if the airport gets at most 1 inch,
    \$862 at 1 to 2 inches, and almost nothing above. \emph{Challenge:}
    today's order book fills only 13\% of it.
  \item \textbf{Advisor.} A farmer types one sentence (``150 acres of hay near
    Weatherford, TX''), and Claude uses our tools to locate the land, check
    the nearest station, and design and explain a contract. Tested live.
\end{itemize}

% ---------------------------------------------------------------------------
\head{2. Are the final results ready?}
Almost, but not yet. We have enough for the presentation: the payouts and the
app work end to end. A formal report needs more research on what these
contracts are worth to real farms and how they compare with subsidized
insurance. Our plan:

{\small
\begin{tabularx}{\textwidth}{@{}lY@{}}
\toprule
Owner & Research for the final report (due 12/15) \\
\midrule
Jack & Model real-world farms accurately and backtest them: what each would
  pay for this insurance, and what it would have paid out in past years \\
Maya & Formalize the gap between our insurance and subsidized PRF: how much
  of it is the subsidy, how much is basis risk, and how much is liquidity \\
Evan & Clean up the app and research more exotic contract ideas, such as
  two-month and combined drought-and-flood contracts \\
\bottomrule
\end{tabularx}
}
\par\vspace{1pt}
All three of us will build the slides for 11/17 and write the final report.

% ---------------------------------------------------------------------------
\head{3. Refinements to our plan}
\begin{itemize}
  \item \textbf{Scope:} instead of copying PRF, we offer drought-weighted
    contracts that cover what PRF leaves open. Liquidity, not price, is now
    the main limit.
  \item \textbf{Responsibilities:} updated to the roles in Section 2.
\end{itemize}

\end{document}
